%%%%%%%%%%%%%%%%%%%%%%%%%%%%%%%%%%%%%%%%%%%%%%%%%%%%%%%
% % Lines starting with % are comments, which are ignored.
% % This is a handy way of indicating the date and version of
% % your document, to wit:
% %
% % LaTeX sample file
% % Modified March, 2002
% %
%%%%%%%%%%%%%%%%%%%%%%%%%%%%%%%%%%%%%%%%%%%%%%%%%%%%%%%
\DocumentMetadata{
  % PDF standards: ONE key whose value is a comma-separated list --
  % UA-2 (accessibility) plus A-4f (archival; the archival companion of
  % UA-2). Do NOT use separate pdfstandard=... keys: with ordinary
  % key-value parsing only the last one takes effect.
  pdfstandard={UA-2,A-4f},
  lang=en-US,
  tagging=on,
  tagging-setup={
    math/setup=mathml-SE
  }
}

\documentclass{ltx-talk}
% lecture-macros.tex -- shared package setup for lecture-template.tex
% (ltx-talk slides / handouts).
%
% NOTE ON PATHS: every template in this repository is stored one level up
% from where it is compiled (see README.md). All relative \input paths
% below are resolved against the folder the document is compiled in,
% which sits one directory above this file.

% --- Tables, bookmarks, strikethrough ------------------------------------
% booktabs: professional table rules (\toprule, \midrule, \bottomrule)
% bookmark: PDF outline entries (PDF bookmarks)
% ulem[normalem]: \sout strikethrough that survives math mode
\usepackage{booktabs}
\usepackage{bookmark}
\usepackage[normalem]{ulem}
\usepackage{multicol}

% --- Math ------------------------------------------------------------
% amsmath: equation environments and numbering
% mathrsfs: \mathfrak  |  stmaryrd: extra symbols
\usepackage{amsmath}
\usepackage{mathrsfs}
\usepackage{stmaryrd}

% --- Layout -------------------------------------------------------------
% float: [H] float placement
% fontspec: Unicode/OpenType fonts (compile with XeLaTeX or LuaLaTeX)
\usepackage{float}
\usepackage{fontspec}

% --- Colors ---------------------------------------------------------------
% dvipsnames: extended named color palette (the Accessible* colors below
% are built on it);
% table: \rowcolors / \cellcolor for shading tables.
% NOTE: xcolor may only be loaded ONCE per document -- a second
% \usepackage{xcolor} with different options is a fatal "Option clash"
% error, so ALL of its options are listed here.
\usepackage[dvipsnames,table]{xcolor}

% --- Multimedia -----------------------------------------------------------
% animate: embeds audio/video (and legacy Flash) media in a frame.
\usepackage{animate}

% --- Bibliography ----------------------------------------------------------
% biblatex (biber backend): citations for lecture notes.
\usepackage[
  backend=biber,
  style=alphabetic,
  url=true,
  doi=true,
  isbn=true,
  backref=true
]{biblatex}

% --- Colored boxes ----------------------------------------------------------
% talkthemetcolorbox: the colored-box theme for ltx-talk slides
% (shadow/rounded/shaded blocks, showtitle, titlepage styles).
\usepackage[
  shadow,
  rounded,
  shaded,
  blocks,
  showtitle,
  titlepage
]{talkthemetcolorbox}

% --- "theindex" helper environment ------------------------------------------
% Provides a lightweight theindex environment (an itemize) and the
% \@idxitem/\subitem/\subsubitem/\indexspace macros that LaTeX's standard
% theindex machinery expects, so hand-written index pages indent correctly.
\makeatletter
\newenvironment{theindex}
  {\begin{itemize}}
  {\end{itemize}}
\providecommand\@idxitem{\par\hangindent 40\p@}
\providecommand\subitem{\@idxitem \hspace*{20\p@}}
\providecommand\subsubitem{\@idxitem \hspace*{30\p@}}
\providecommand\indexspace{\par \vskip 10\p@ \@plus5\p@ \@minus3\p@\relax}
\makeatother

% --- Index -----------------------------------------------------------------
% imakeidx: builds the printed index.
% The redefinition below appends "|hyperpage" to every index entry so that,
% in the PDF, each index entry is a link to its page.
\usepackage{imakeidx}
\makeatletter
\let\imki@orig@wrindexentrysplit\imki@wrindexentrysplit
\renewcommand\imki@wrindexentrysplit[3]{%
  \imki@orig@wrindexentrysplit{#1}{#2|hyperpage}{#3}%
}

\makeatletter
% --- Section titles: capture + one bookmark per logical section -------------
% \section is redefined so that:
%   1. the section title is captured into the plain-TeX register
%      \mytalksectitletoks (reprinted verbatim on the auto divider frame,
%      see the cmd/section/after hook below -- plain-TeX storage is
%      catcode-safe, unlike expanding the title in a hook), and
%   2. exactly ONE depth-1 PDF bookmark is created per \section call
%      (hyperref's own bookmark is suppressed via \Hy@writebookmark).
\DeclareCommandCopy{\OriginalSectionForBookmarks}{\section}
\newtoks\mytalksectitletoks

\RenewDocumentCommand{\section}{s o m}{%
  \begingroup
    \def\Hy@writebookmark##1##2##3##4##5{}%
    \global\mytalksectitletoks={#3}%
    \IfBooleanTF{#1}
      {\phantomsection\IfNoValueTF{#2}{\OriginalSectionForBookmarks*{#3}}{\OriginalSectionForBookmarks*[#2]{#3}}}%
      {\IfNoValueTF{#2}{\OriginalSectionForBookmarks{#3}}{\OriginalSectionForBookmarks[#2]{#3}}}%
  \endgroup
  \IfBooleanTF{#1}
    {\IfNoValueTF{#2}{\pdfbookmark[1]{#3}{section-star.\thepage.\number\inputlineno}}{\pdfbookmark[1]{#2}{section-star.\thepage.\number\inputlineno}}}%
    {\IfNoValueTF{#2}{\pdfbookmark[1]{Section \thesection: #3}{section.\thesection}}{\pdfbookmark[1]{Section \thesection: #2}{section.\thesection}}}%
}
\makeatother

% --- One PDF bookmark per logical frame (NOT per overlay) ---------------------
% ltx-talk writes one PDF page per overlay, so a naive frametitle hook would
% register a new bookmark for every overlay of the same frame. The
% \g_slideutils_frame_bookmark_int register remembers the last frame that
% got a bookmark, so subsequent overlays of the same frame are skipped.
% This is also what keeps BeamerPresenter's cumulative annotations working:
% one logical frame must stay one logical unit in the PDF outline.
\ExplSyntaxOn
\int_new:N \g_slideutils_frame_bookmark_int

\AddToHook{cmd/frametitle/after}{%
  \tl_if_empty:NF \g__talk_frame_title_tl
    {
      \ifnum\value{frame}=\g_slideutils_frame_bookmark_int\relax
      \else
        \int_gset:Nn \g_slideutils_frame_bookmark_int { \value{frame} }%
        \pdfbookmark[2]{\tl_use:N \g__talk_frame_title_tl}
          {frame.\int_use:N \g_slideutils_frame_bookmark_int}
      \fi
    }
}
\ExplSyntaxOff
\makeatother

% --- Printed index + theme for the colored boxes -------------------------------
% \makeindex[...]: two-column index titled "Index", listed in the TOC.
% The two EditInstance calls style the standard header/footer instances
% of talkthemetcolorbox (cyan background; footer order title/author/
% date/framenumber).
\makeindex[columns=2, title=Index, intoc]\EditInstance{header}{std}{
  background-color = cyan,
  color            = black,
  height           = 0.75cm
}
\EditInstance{footer}{std}{
  background-color = cyan,
  element-order = {title, author, date, framenumber},
}

% --- Auto section divider frame ---------------------------------------------------
% After each \section, insert a centered, large, structure-colored frame
% that prints the section title (from \mytalksectitletoks).
\AddToHook{cmd/section/after}{%
  \begin{frame}
    \centering
    \color{structure}
    \LARGE
    \the\mytalksectitletoks
  \end{frame}%
}

% --- Tagging sockets (tagged PDF structure) ---------------------------------------
% ltx-talk exposes tagging "sockets" where the PDF structure can attach
% section titles; register one and leave it unassigned (none) so the
% class's own tagging keeps control.
\ExplSyntaxOn
\AtBeginDocument{
  \NewTaggingSocketPlug { talk / sec / title } { none } { }
  \AssignTaggingSocketPlug { talk / sec / title } { none }
}
\ExplSyntaxOff

% --- Accessibility roles ----------------------------------------------------------------
% Map \frametitle to the H1 heading role in the tagged PDF, so screen
% readers treat slide titles as headings.
\tagpdfsetup{role / new-tag = frametitle / H1}

\makeatletter
% --- Page label = frame number -------------------------------------------------------
% Redefine what feeds the PDF page label to track the frame counter,
% not the physical page counter. Replace \theframenumber below with
% whatever internal macro/property ltx-talk actually uses to produce
% the "framenumber" footer element — confirm via ltx-talk-frame.dtx.
\renewcommand\thepage{\arabic{frame}}
\makeatother

\addbibresource{refs.bib}

% --- "Problem Time!" slides ------------------------------------------------------------------
% \ProblemSlide{problem}{solution} emits two frames: the problem, then
% its solution after a \pause, each numbered by the problem counter.
\newcounter{problemcounter}
\setcounter{problemcounter}{0}
\newcommand{\ProblemSlide}[2]{
  \stepcounter{problemcounter}
  \begin{frame}
    \frametitle{Problem Time!}
    Problem \theproblemcounter: #1
  \end{frame}

  \begin{frame}
    \frametitle{Problem Time!}
    Solution to Problem \theproblemcounter: \pause

    #2
  \end{frame}
}
% \highlight[<colour>]{<stuff>}
% from https://tex.stackexchange.com/questions/33401/a-version-of-colorbox-that-works-inside-math-environments
\newcommand{\highlight}[2][yellow]{\mathchoice%
  {\colorbox{#1}{$\displaystyle#2$}}%
  {\colorbox{#1}{$\textstyle#2$}}%
  {\colorbox{#1}{$\scriptstyle#2$}}%
  {\colorbox{#1}{$\scriptscriptstyle#2$}}}%

% --- Fonts ------------------------------------------------------------------------------
% Lato (text) + LeteSansMath (math): open, accessible typefaces.
\setmainfont{Lato}[
  Ligatures=TeX,
  Numbers={Lining,Monospaced}
]
\setsansfont{Lato}[
  Ligatures=TeX,
  Numbers={Lining,Monospaced}
]
\setmathfont{LeteSansMath.otf}\usepackage{emoji}

% --- Accessible colors ---------------------------------------------------------------------
% WCAG-contrast-safe variants of red/green for text on white:
% red!80!black and green!40!black both clear the 4.5:1 requirement.
% AltYellow: soft background for highlight boxes.
\colorlet{AccessibleRed}{red!80!black}
\colorlet{AccessibleGreen}{green!40!black}
\definecolor{AltYellow}{RGB}{255,245,170}

% --- Links ---------------------------------------------------------------------------------
% hyperref + link colors: every link uses AccessibleGreen (7.2:1 on
% white), so internal, citation, and URL links all pass WCAG 1.4.3.
\usepackage{hyperref}
\hypersetup{
  colorlinks=true,
  linkcolor=AccessibleGreen,
  urlcolor=AccessibleGreen,
  citecolor=AccessibleGreen,
}

% --- Shared course macros ---------------------------------------------------------------------
% Root macros.tex: math shorthands used across all course documents
% (\RR, \NN, \ZZ, \QQ, \grad, \curl, \tr, \der, \pder, ...).
\newcommand{\RR}{\mathbb{R}}
\newcommand{\NN}{\mathbb{N}}
\newcommand{\ZZ}{\mathbb{Z}}
\newcommand{\QQ}{\mathbb{Q}}
\newcommand{\vv}{\vec{v}}
\newcommand{\ww}{\vec{w}}
\newcommand{\xx}{\vec{x}}
\newcommand{\yy}{\vec{y}}
\newcommand{\zz}{\vec{z}}
\newcommand{\uu}{\vec{u}}
\newcommand{\basis}[1]{\vec{e_{#1}}}
\newcommand{\id}{\mathrm{id}}
\newcommand{\CC}{\mathbb{C}}
\renewcommand{\div}{\operatorname{div}}
\newcommand{\grad}{\operatorname{grad}}
\newcommand{\curl}{\operatorname{curl}}
\newcommand{\tr}{\operatorname{tr}}
\newcommand{\der}[2]{\dfrac{d #1}{d #2}}
\newcommand{\pder}[3][]{\dfrac{\partial^{#1} #2}{\partial #3}}


% % Title and author(s)
%%%%%%%%%%%%%%%%%%%%%%%%%%%%%%%%%%%%%%%%%%%%%%%%%%%%%%%
\title{Lecture 14 - L'Hôpital's Rule}
\author{Garrett Credi}
\date{10/27/2026}
%%%%%%%%%%%%%%%%%%%%%%%%%%%%%%%%%%%%%%%%%%%%%%%%%%%%%%%
% %
% % The next command allows your in import encapsulated
% % postscript files, .epsf or .eps files, which
% % contain vector graphic image data.
% %
%%%%%%%%%%%%%%%%%%%%%%%%%%%%%%%%%%%%%%%%%%%%%%%%%%%%%%%
% %   The first thanks indicates your affiliation
% %
% %  Just the name here.
% %
% % Your mailing address goes at the end.
% %
% % \thanks is also how you indicate grant support
% %
%%%%%%%%%%%%%%%%%%%%%%%%%%%%%%%%%%%%%%%%%%%%%%%%%%%%%%%

%%%%%%%%%%%%%%%%%%%%%%%%%%%%%%%%%%%%%%%%%%%%%%%%%%%%%%%
% $ The next command ensures the margins are 1 inch on
% % all sides.
% %
%%%%%%%%%%%%%%%%%%%%%%%%%%%%%%%%%%%%%%%%%%%%%%%%%%%%%%%
\begin{document}
\begin{frame}
  \maketitle
\end{frame}

\begin{frame}
  \frametitle{Updates}

  We have Exam 2 next week!

  Practice exams and the review sheets are both out on Canvas,
  compiled lecture will be uploaded after class.

  Solutions to the practice exam and review sheet will be uploaded Thursday night 
  to give you all time to try out the problems before looking at the solutions.

  If you're feeling worried about the exam, please feel free to come to my office hours or email me to set up a time to meet :)

  Re: quizzes. I'll grind out quiz grading tonight and tomorrow. Let me know if you want me to give feedback on a specific quiz of yours earlier though :)
\end{frame}

\begin{frame}
  \frametitle{Quiz Time}

  Let's get into neat rows, spread out, and slay this quiz :)
\end{frame}

\section{L'Hôpital's Rule}

\begin{frame}
  \frametitle{Limit Review I}

  Recall from the first two weeks of class:

  we use \textit{limits} to determine the values
  that functions approximate as their inputs approximate a given value. \pause

  They were very useful in examining \textit{indeterminate forms}:
  expressions like $\dfrac{f(x)}{g(x)}$ around values $x = x_0$
  where $f(x_0) = g(x_0) = 0$ or $\infty$. \pause

  Our tools included:
  \begin{multicols}{2}
    \begin{itemize}
      \item Factoring out $(x - x_0)$
      \item Replacing the limit with the limit of a continuous function
      \item Addition/subtraction/product/quotient computations
      \item Squeeze theorem
      \item Etc.
    \end{itemize}
  \end{multicols}
\end{frame}

\begin{frame}
  \frametitle{Limit Review II}

  We also defined derivatives in terms of limits: \pause
  \begin{equation}
    f'(x) = \lim_{h \rightarrow 0} \dfrac{f(x + h) - f(x)}{h}
  \end{equation}\pause

  But then quickly pivoted to constructing and working with the \textit{synthetic} derivative rules, e.g.: \pause
  \begin{multicols}{2}
    \begin{itemize}
      \item Power Rule: $\dfrac{d}{dx} x^n = nx^{n-1}$
      \item Exponential/Logarithm rule: $\dfrac{d}{dx} e^x = e^x$, $\dfrac{d}{dx} \ln(x) = \dfrac{1}{x}$
      \item Trig Rules: $\dfrac{d}{dx} \sin(x) = \cos(x)$
      \item Sum/Difference rule: $\dfrac{d}{dx}(f(x) \pm g(x)) = \dfrac{d}{dx}f(x) \pm \dfrac{d}{dx}g(x)$
      \item Product Rule: $\dfrac{d}{dx} f(x)g(x) = \left(\dfrac{d}{dx} f(x)\right) \cdot g(x) + f(x) \cdot \left(\dfrac{d}{dx}g(x)\right)$
      \item Chain Rule: $\dfrac{d}{dx} g(f(x)) = \dfrac{dg}{du} \mid_{u = f(x)} \cdot \dfrac{df}{dx}$.
    \end{itemize}
  \end{multicols}
\end{frame}

\begin{frame}
  \frametitle{Derivative Review I}

  The importance of the derivative came from it's interpretation as
  \textit{the best linear approximation}: \pause
  \begin{equation}
    \lim\limits_{h \rightarrow 0} \dfrac{|f(x) - (f(x_0) + f'(x_0)(x - x_0))|}{h} = 0
  \end{equation} \pause

  or, in practical terms:
  \begin{equation}
    f(x) = f(x_0) + f'(x_0)(x - x_0) \pause + O((x - x_0)^2)
  \end{equation} \pause

  We also saw the \textit{quadratic approximation}:
  \begin{equation}
    f(x) = f(x_0) + f'(x_0)(x - x_0) + \dfrac{f''(x_0)}{2} (x - x_0)^2 \pause + O((x - x_0)^3)
  \end{equation}
\end{frame}

\begin{frame}
  \frametitle{L'Hôpital's Rule I}

  Because of our ability to approximate functions via their linear/quadratic approximations, we gain another tool to put in our limit computing toolkit
  when we encounter indeterminate forms. \pause

  If $f(x)$ and $g(x)$ are two functions that are: \pause
  \begin{enumerate}
    \item $0$ when $x = x_0$ \pause
    \item Differentiable at $x = x_0$ \pause
  \end{enumerate}

  then, assuming $g'(x_0) \neq 0$, an application of squeeze theorem
  gives us:
  \begin{equation}
    \lim\limits_{x \rightarrow x_0} \dfrac{f(x)}{g(x)} = \lim\limits_{x \rightarrow x_0} \dfrac{f(x_0) + f'(x_0)(x - x_0)}{g(x_0) + g'(x_0)(x - x_0)} \pause = \lim\limits_{x \rightarrow x_0} \dfrac{f'(x_0)(x - x_0)}{g'(x_0)(x - x_0)} = \lim\limits_{x \rightarrow x_0} \dfrac{f'(x_0)}{g'(x_0)} 
  \end{equation} \pause

  We can also work around $g'(x_0)$ being zero by instead reducing
  \begin{equation}
    \lim\limits_{x \rightarrow x_0} \dfrac{f(x)}{g(x)} = \lim\limits_{x \rightarrow x_0} \dfrac{f'(x)}{g'(x)}
  \end{equation}
\end{frame}

\begin{frame}
  \frametitle{L'Hôpital's Rule II}

  \begin{alertblock}{(Mini) L'Hôpital's Rule}\index{L'Hôpital's Rule!Mini}
    If $f(x)$, $g(x)$ are two functions that are:
    \begin{enumerate}
      \item Zero at $x = x_0$ (i.e. $f(x_0) = g(x_0) = 0$)
      \item Differentiable at $x = x_0$
    \end{enumerate}
    then:
    \begin{equation}
      \lim\limits_{x \rightarrow x_0} \dfrac{f(x)}{g(x)} = \lim\limits_{x \rightarrow x_0} \dfrac{f'(x)}{g'(x)}
    \end{equation}
  \end{alertblock} \pause

  In other words, we can compute limits of indeterminate forms by \textit{differentiating the two functions} to kill off any obstructions to continuity.
\end{frame}

\begin{frame}
  \frametitle{L'Hôpital's Rule III}

  This comes from a generalization of both linear and quadratic approximation,
  arguably the most important concept taught in Calculus II:

  \begin{alertblock}{Taylor's Theorem}\index{Taylor's Theorem}\index{Taylor Polynomial}
    For $f(x)$ a function whose higher derivatives $\dfrac{d^n f}{dx^n} \mid_{x = x_0}$ all exist one can form the \textit{Taylor Polynomial} of $f(x)$ at $x = x_0$ defined as:
    \begin{equation}
      T_f(x) := \sum\limits_{n = 0}^{\infty} \dfrac{f^{(n)}(x_0)}{n!} (x - x_0)^n
    \end{equation}
  and as $x \rightarrow x_0$ the error term $E(x) = |f(x) - T_f(x)|$
  approaches zero \textit{extremely rapidly}. (And in most nice cases it turns out that $f(x) = T_f(x)$ \textit{literally} on a small interval containing $x_0$).
  \end{alertblock}

  (this is not material you are expected to remember for this class)
\end{frame}

\begin{frame}
  \frametitle{L'Hôpital's Rule IV}

  L'Hôpital's Rule is then just what happens when you compute $\lim\limits_{x \rightarrow x_0} \dfrac{f(x)}{g(x)}$
  by replacing it with $\lim\limits_{x \rightarrow x_0} \dfrac{T_f(x)}{T_g(x)}$ using the Squeeze Theorem.
  E.g.:
  \begin{align}
    \lim\limits_{x \rightarrow x_0} \dfrac{f(x)}{g(x)} &= \lim\limits_{x \rightarrow x_0} \dfrac{f(x_0) + f'(x_0)(x - x_0) + \dfrac{f''(x_0)}{2}(x - x_0)^2 + \dfrac{f'''(x_0)}{6}(x - x_0)^3 + \ldots}{g(x_0) + g'(x_0)(x - x_0) + \dfrac{g''(x_0)}{2} (x - x_0)^2 + \dfrac{g'''(x_0)}{6}(x - x_0)^3 + \ldots} \\ \pause
                                                       &= \lim\limits_{x \rightarrow x_0} \dfrac{f'(x_0)(x - x_0) + \dfrac{f''(x_0)}{2}(x - x_0)^2 + \dfrac{f'''(x_0)}{6}(x - x_0)^3 + \ldots}{g'(x_0)(x - x_0) + \dfrac{g''(x_0)}{2}(x - x_0)^2 + \dfrac{g'''(x_0)}{6}(x - x_0)^3 + \ldots} \\ \pause
                                                       &= \lim\limits_{x \rightarrow x_0} \dfrac{f'(x_0) + \dfrac{f''(x_0)}{2} (x - x_0) + \dfrac{f'''(x_0)}{6} (x - x_0)^2 + \ldots}{g'(x_0) + \dfrac{g''(x_0)}{2} (x - x_0) + \dfrac{g'''(x_0)}{6} (x - x_0)^2 + \ldots}
  \end{align}
\end{frame}
\begin{frame}
  \frametitle{L'Hôpital's Rule V}
  What if $\lim\limits_{x \rightarrow x_0} \dfrac{f'(x)}{g'(x)}$ is another $\dfrac{0}{0}$ indeterminate form? \pause e.g. $f'(x_0) = g'(x_0) = 0$? \pause

  \begin{align}
    \lim\limits_{x \rightarrow x_0} \dfrac{f(x)}{g(x)} &= \lim\limits_{x \rightarrow x_0} \dfrac{f'(x_0) + \dfrac{f''(x_0)}{2}(x - x_0) + \dfrac{f'''(x_0)}{6}(x - x_0)^2 + \ldots}{g'(x_0) + \dfrac{g''(x_0)}{2}(x - x_0) + \dfrac{g'''(x_0)}{6}(x - x_0)^2 + \ldots} \\ \pause
                                                       &= \lim\limits_{x \rightarrow x_0} \dfrac{\dfrac{f''(x_0)}{2}(x - x_0) + \dfrac{f'''(x_0)}{6}(x - x_0)^2 + \ldots}{\dfrac{g''(x_0)}{2}(x - x_0) + \dfrac{g'''(x_0)}{6}(x - x_0)^2 + \ldots} \\ \pause
                                                       &= \lim\limits_{x \rightarrow x_0}{f''(x_0) + \dfrac{f'''(x_0)}{3}(x - x_0) + \ldots}{g''(x_0) + \dfrac{g'''(x_0)}{3}(x - x_0) + \ldots} \pause
  \end{align}

  In other words, we just L'Hôpital it again!
\end{frame}

\begin{frame}
  \frametitle{L'Hôpital's Rule VI}

  This gives us a general procedure to compute limits of the form $\lim_{x \rightarrow x_0}\dfrac{f(x)}{g(x)}$: \pause

  \begin{alertblock}
    {L'Hôpital's Process for Computing Limits of the Indeterminate Form $\dfrac{0}{0}$}\index{L'Hôpital's Rule!Process for 0/0 forms}
    To compute $\lim_{x \rightarrow x_0} \dfrac{f(x)}{g(x)}$:
    \begin{enumerate}
      \item Compute $f(x_0)$ and $g(x_0)$. \pause
      \item If $f(x_0) \neq 0$ and $g(x_0) \neq 0$, \pause then $\lim_{x \rightarrow x_0} \dfrac{f(x)}{g(x)} = \dfrac{f(x_0)}{g(x_0)}$. \pause
      \item If $f(x_0) \neq 0$ but $g(x_0) = 0$ then we're ending up with something like $\dfrac{N}{0}$ so: \pause $\lim_{x \rightarrow x_0}\dfrac{f(x)}{g(x)}$ DNE \pause
      \item If $f(x_0) = g(x_0) = 0$ and $f(x)$, $g(x)$ are differentiable at $x = x_0$ \pause, replace $f(x)$ with $f'(x)$ and $g(x)$ with $g'(x)$
        and restart step 1.
    \end{enumerate} \pause
  \end{alertblock}
\end{frame}

\begin{frame}
  \frametitle{L'Hôpital's Rule VII}

  What about other indeterminate forms like: \pause
  \begin{equation*}
    \dfrac{\infty}{\infty}, \quad 0 \cdot \infty, \quad 0^0, \quad \infty^0, \quad 1^\infty, \quad \infty - \infty?
  \end{equation*} \pause

  If L'Hôpital's rule helps us compute limits of indeterminate forms like $\dfrac{0}{0}$ \pause why not just try and force those above forms into something related enough and that looks like $\dfrac{0}{0}$! \pause

  E.g. if we had $\lim\limits_{x \rightarrow 0} \dfrac{\cot(x)}{\ln(x^2)}$ \pause
  we see this is of the form $\dfrac{\infty}{\infty}$ \pause but we can change it to:
  \begin{equation*}
    \lim\limits_{x \rightarrow 0} \dfrac{\cot(x)}{\ln(x^2)} = \left(\lim\limits_{x \rightarrow 0} \dfrac{\frac{1}{\ln(x^2)}}{\frac{1}{\cot(x)}}\right)
  \end{equation*} \pause
  which is of the form $\dfrac{0}{0}$.
\end{frame}

\begin{frame}
  \frametitle{L'Hôpital's Rule Examples I}
  We can then compute:
  \begin{align}
    \lim\limits_{x \rightarrow 0} \dfrac{\cot(x)}{\ln(x^2)} &= \lim\limits_{x \rightarrow 0} \dfrac{\frac{1}{\ln(x^2)}}{\tan(x)} \\ \pause
                                                            &= \lim\limits_{x \rightarrow 0} \dfrac{\dfrac{d}{dx}(\ln(x^2))^{-1}}{\dfrac{d}{dx} \tan(x)} \\ \pause
                                                            &= \lim\limits_{x \rightarrow 0} \dfrac{2x \cdot \dfrac{1}{(x^2)} \cdot -1 \cdot (\ln(x^2))^{-2}}{\sec(x)^2} \\ \pause
                                                            &= \lim\limits_{x \rightarrow 0}\dfrac{ \dfrac{-2 \ln(x^2)}{x}}{\sec(x)^2}
  \end{align}
  since $\sec(0)^2 = 1$, and since the numerator is an indeterminate form of type $\dfrac{\infty}{0}$, we can conclude this limit DNE.
\end{frame}

\begin{frame}
  \frametitle{L'Hôpital's Rule Examples II}
  If you're a fan of Mean Girls, you might be interested in computing:
  \begin{equation}
    \lim\limits_{x \rightarrow 0} \dfrac{\ln(1 - x) - \sin(x)}{1 - \cos(x)^2}
  \end{equation} \pause

  We check @ $x = 0$, we'd get:
    $\dfrac{\ln(1 - 0) - \sin(0)}{1 - \cos(0)^2} = \pause \dfrac{\ln(1) - 0}{1 - 1^2} = \pause \dfrac{0}{0}$
  so this is already in a form amenable to L'Hôpital's Rule! \pause
  \begin{align}
    \lim\limits_{x \rightarrow 0} \dfrac{\ln(1 - x) - \sin(x)}{1 - \cos(x)^2} &= \lim\limits_{x \rightarrow 0} \dfrac{\dfrac{d}{dx} (\ln(1 - x) - \sin(x))}{\dfrac{d}{dx} (1 - \cos(x)^2)} \\ \pause
                                                                                            &= \lim\limits_{x \rightarrow 0} \dfrac{\frac{-1}{1 - x} - \cos(x)}{-2 \cos(x) \cdot - \sin(x)}
  \end{align} \pause
  Which at $0$ is $\dfrac{\dfrac{-1}{1 - 0} - \cos(0)}{2 \cos(0) \sin(0)} = \dfrac{-2}{0}$ so \pause the limit does not exist!
\end{frame}

\begin{frame}
  \frametitle{L'Hôpital's Rule Examples III}
  Another one:
  \begin{equation}
    \lim\limits_{x \rightarrow 0} x^x
  \end{equation} \pause

  Indeterminate form of type \pause $0^0$. \pause What can we use to turn exponentiation into a product/quotient of things? \pause Logarithms! \pause

  \begin{align}
    \ln(\lim_{x \rightarrow 0} x^x) &= \lim_{x \rightarrow 0} \ln(x^x) \pause = \lim_{x \rightarrow 0} x \ln(x) \\ \pause
                                    &= \lim_{x \rightarrow 0} \dfrac{\ln(x)}{\frac{1}{x}} \pause = \lim_{x \rightarrow 0} \dfrac{\frac{d}{dx} \ln(x)}{\frac{d}{dx} \frac{1}{x}} \pause \\
                                    &= \lim_{x \rightarrow 0} \dfrac{\frac{1}{x}}{-\frac{1}{x^2}} \pause = \lim_{x \rightarrow 0}-\frac{x^2}{x} = 0
  \end{align} \pause

  So $\ln(\lim_{x \rightarrow 0} x^x) = 0$, which means \pause $\lim_{x \rightarrow 0} x^x = e^0 = 1$.
\end{frame}

\begin{frame}
  \frametitle{L'Hôpital's Rule Cheat Method!}
  In the previous example we did literal L'Hôpital's rule
  on an indeterminate form of type $\dfrac{\infty}{\infty}$. Why is this valid? \pause

  \begin{align}
    \lim_{x \rightarrow x_0} \dfrac{f(x)}{g(x)} &= \lim_{x \rightarrow x_0} \dfrac{\frac{1}{g(x)}}{\frac{1}{f(x)}} \\ \pause
                                                &= \lim_{x \rightarrow x_0}
                                                \dfrac{\frac{d}{dx} \frac{1}{g(x)}}{\frac{d}{dx}\frac{1}{f(x)}} \\ \pause
                                                &= \lim_{x \rightarrow x_0} \dfrac{\frac{-g'(x)}{g(x)^2}}{\frac{-f'(x)}{f(x)^2}} = \left(\lim_{x \rightarrow x_0} \dfrac{g'(x)}{f'(x)}\right) \cdot \left(\lim_{x \rightarrow x_0} \dfrac{f(x)}{g(x)}\right)^2
  \end{align}
  from which you can solve that, for $\lim_{x \rightarrow x_0} f(x) = \lim_{x \rightarrow x_0} g(x) = \infty$
  \begin{equation}
    \lim_{x \rightarrow x_0} \dfrac{f(x)}{g(x)} = \lim_{x \rightarrow x_0} \dfrac{f'(x)}{g'(x)}
  \end{equation}
\end{frame}

\begin{frame}
  \frametitle{Superpowered L'Hôpital's Rule}

  \begin{alertblock}{L'Hôpital's Rule}\index{L'Hôpital's Rule}
    For $f(x)$, $g(x)$ two functions where \textit{either}:
    \begin{enumerate}
      \item $\lim_{x \rightarrow x_0} f(x) = \lim_{x \rightarrow x_0} g(x) = 0$ and $f(x), g(x)$ are differentiable at $x = x_0$
      \item $\lim_{x \rightarrow x_0} |f(x)| = \lim_{x \rightarrow x_0} |g(x)| = \infty$ and $\frac{1}{f(x)}, \frac{1}{g(x)}$ are differentiable at $x = x_0$.
    \end{enumerate}
    Then we have that:
    \begin{equation}
      \lim_{x \rightarrow x_0} \dfrac{f(x)}{g(x)} = \lim_{x \rightarrow x_0} \dfrac{f'(x)}{g'(x)}
    \end{equation}
  \end{alertblock} \pause

  Then to compute the limit of \textit{any} indeterminate form, we need only
  equate it to some indeterminate form of type $\dfrac{0}{0}$ or $\dfrac{\infty}{\infty}$ and use L'Hôpital's Rule to compute it.
\end{frame}

\begin{frame}
  \frametitle{L'Hôpital's Rule Examples IV}
  Let's consider $\lim_{x \rightarrow 0} \dfrac{\sin(x)}{\cos(x)}$
  as well as the following "computation": \pause

  \begin{align}
    \lim\limits_{x \rightarrow x_0} \dfrac{\sin(x)}{\cos(x)} &= \lim\limits_{x \rightarrow x_0} \dfrac{\frac{d}{dx} \sin(x)}{\frac{d}{dx} \cos(x)} \\ \pause
                                                             &= \lim\limits_{x \rightarrow 0} \dfrac{\cos(x)}{-\sin(x)}
  \end{align}
  DNE since $\cos(0) = 1$ while $-\sin(0) = 0$ so we end up with $\frac{1}{0}$ so DNE. Thus $\tan(0)$ does not exist. \pause

  \textcolor{red}{WE DIDN'T ACTUALLY CHECK IF $\dfrac{\sin(x)}{\cos(x)}$ WAS AN INDETERMINATE FORM AT $x = 0$!} \pause

  $\sin(0) = 0$ and $\cos(0) = 1$ so $\dfrac{\sin(x)}{\cos(x)}$ is \textit{continuous} at $x = 0$ so $\lim_{x \rightarrow 0} \dfrac{\sin(x)}{\cos(x)} = \dfrac{\sin(0)}{\cos(0)} = \dfrac{0}{1} = 0 = \tan(0)$.
\end{frame}

\begin{frame}
  \frametitle{L'Hôpital's Rule Examples V}
  Finally, let's consider
  \begin{equation}
    \lim_{x \rightarrow \infty} \left(1 + \dfrac{1}{x}\right)^x
  \end{equation}
  (which actually matters since its the ``rate of continuously compounding interest") \pause
  
  What type of indeterminate form? \pause $\lim_{x \rightarrow \infty} 1 + \frac{1}{x} = 1$, so \pause $1^\infty$.

  How do we turn exponentiation into a product/quotient? \pause Log!

  \begin{align}
    \ln(\lim_{x \rightarrow \infty} \left(1 + \dfrac{1}{x}\right)^x) &= \lim_{x \rightarrow \infty} \ln(\left(1 + \dfrac{1}{x}\right)^x) \\ \pause
                                                                     &= \lim_{x \rightarrow \infty} x \ln(1 + \frac{1}{x}) \pause = \lim_{x \rightarrow \infty} \dfrac{\ln(1 + \frac{1}{x})}{\frac{1}{x}} \\ \pause
                                                                     &= \lim\limits_{x \rightarrow \infty} \dfrac{-\frac{1}{x^2} \cdot \frac{1}{1 + \frac{1}{x}}}{-\frac{1}{x^2}} \pause = \lim_{x \rightarrow \infty} \dfrac{x}{x + 1} = \pause 1
  \end{align} \pause

  So $\ln\left(\lim_{x \rightarrow \infty} \left( 1 + \dfrac{1}{x}\right)^x \right) = 1$ and therefore \pause $\lim_{x \rightarrow \infty} \left(1 + \dfrac{1}{x}\right)^x = e$.
\end{frame}

\begin{frame}
  \frametitle{Worksheet Time!}

  Let's get into groups and work on the worksheet :)
\end{frame}

\begin{frame}
  \frametitle{End of Class}

  That's all folks :)
\end{frame}
\end{document}
